% ============================================================
% Come si usano, argomento 22: stima classica e limite di Cramér-Rao
% ============================================================
\argomento{Stima classica e limite di Cramér-Rao}{Formule a p.~\pageref*{f:verosimiglianza}}

\begin{scelta}{Come scegliere la formula}
\renewcommand{\arraystretch}{1.25}
\begin{tabular}{@{}P{0.55\linewidth}P{0.39\linewidth}@{}}
  \toprule
  \textbf{Se il testo\ldots} & \textbf{si usa\ldots} \\
  \midrule
  dice che il parametro non è aleatorio e chiede la stima più plausibile & \vv{verosimiglianza} \vv{stimaml} \\
  chiede se uno stimatore sbaglia in media & la polarizzazione \vv{polarizzazione} \\
  chiede l'errore quadratico medio di uno stimatore & \vv{msebias} \\
  chiede il minimo errore possibile, o se lo stimatore è efficiente & \vv{fisher} \vv{cramerrao} \\
  stima la probabilità di testa o la media di dati gaussiani & \vv{mlbernoulli} \vv{mlgauss} \\
  \bottomrule
\end{tabular}

Lo schema tipico: log-verosimiglianza, derivata prima nulla (la stima), media e varianza della stima, derivata seconda
per l'informazione di Fisher, confronto con il limite di Cramér-Rao.
\end{scelta}

\scheda{verosimiglianza}{
  \usoquando{Il parametro non è aleatorio (non c'è una densità a priori) e si vuole sapere quanto ciascun valore rende
    plausibili i dati osservati.}
  \usocome{Si scrive la densità (o la pmf) dei dati osservati e la si legge come funzione di $\theta$; con dati
    indipendenti conviene il logaritmo.}
  \usoattenzione{Come funzione di $\theta$ non è una densità: in generale non ha area 1.}
  \usoesempio{2 teste in 3 lanci: $L(\beta) = \beta^2 (1 - \beta)$. $L(0{,}5) = 0{,}125$ e $L\left( \frac{2}{3}
    \right) = \frac{4}{27} \simeq 0{,}148$: $\beta = \frac{2}{3}$ rende i dati più plausibili.}
}

\scheda{stimaml}{
  \usoquando{Si deve stimare un parametro non aleatorio: è lo stimatore classico più usato.}
  \usocome{Si prende il logaritmo della verosimiglianza, si deriva rispetto a $\theta$ e si annulla la derivata
    (\vv{argminmax}).}
  \usoesempio{Dati esponenziali $1, 2, 3$ di parametro $\lambda$: $\Lambda = 3 \ln\lambda - \lambda (1 + 2 + 3)$, la
    derivata $\frac{3}{\lambda} - 6$ si annulla in $\hat\lambda = \frac{1}{2}$, l'inverso della media dei dati.}
}

\scheda{polarizzazione}{
  \usoquando{Si chiede se lo stimatore sbaglia in media, cioè se ripetendo l'esperimento la media delle stime sarebbe il
    valore vero.}
  \usocome{Si calcola la media dello stimatore e la si confronta con $\theta$.}
  \usoesempio{$\frac{w}{n}$ ha media $\beta$: non polarizzato. $\frac{w + 1}{n + 2}$ ha media $\frac{n\beta + 1}{n +
    2} = \beta + \frac{1 - 2\beta}{n + 2}$: polarizzato (tranne che con $\beta = \frac{1}{2}$), ma asintoticamente non
    polarizzato.}
}

\scheda{msebias}{
  \usoquando{Si vuole misurare quanto le stime sono lontane dal valore vero, sia per lo sparpagliamento sia per lo
    spostamento della media.}
  \usocome{Si calcolano varianza e polarizzazione e si sommano varianza e quadrato della polarizzazione.}
  \usoesempio{$\frac{w + 1}{n + 2}$ con $n = 10$: varianza $\frac{10\, \beta(1 - \beta)}{144}$ e polarizzazione
    $\frac{1 - 2\beta}{12}$; con $\beta = 0{,}5$ l'errore è $\frac{2{,}5}{144} \simeq 0{,}017$.}
}

\scheda{fisher}{
  \usoquando{Serve per il limite di Cramér-Rao: misura quanta informazione i dati portano sul parametro.}
  \usocome{Di solito è più facile la seconda forma: si deriva due volte la log-verosimiglianza, si cambia segno e si fa
    la media rispetto ai dati.}
  \usoesempio{Dati gaussiani: $\frac{\partial^2}{\partial\mu^2} \log f = -\frac{n}{\sigma^2}$, che non dipende dai dati,
    quindi $I_n(\mu) = \frac{n}{\sigma^2}$. Con $\sigma = 2$ e $n = 100$: $I_n = 25$.}
}

\scheda{cramerrao}{
  \usoquando{Si chiede quanto può essere piccolo l'errore di uno stimatore, o se uno stimatore è efficiente.}
  \usocome{Si confronta la varianza dello stimatore con $\frac{1}{I_n(\theta)}$: se coincidono e lo stimatore è non
    polarizzato, è efficiente.}
  \usoattenzione{$\frac{1}{I_n}$ vale per gli stimatori non polarizzati: uno polarizzato può scendere sotto. Con $n =
    10$ e $\beta = 0{,}5$, $\frac{w + 1}{n + 2}$ ha errore $0{,}017$, meno di $\frac{1}{I_n} = 0{,}025$.}
  \usoesempio{Media di 100 dati gaussiani con $\sigma = 2$: nessuno stimatore non polarizzato ha varianza sotto
    $\frac{1}{25} = 0{,}04$ (deviazione standard $0{,}2$); la media campione la raggiunge.}
}

\scheda{mlbernoulli}{
  \usoquando{Si stima una probabilità di successo non aleatoria con la frequenza dei successi.}
  \usocome{È la stessa derivata della stima MAP della moneta misteriosa; la varianza viene dalla binomiale.}
  \usoesempio{30 teste in 100 lanci: $\hat\beta = 0{,}3$; con $\beta \simeq 0{,}3$ la varianza è $\frac{0{,}3 \cdot
    0{,}7}{100} = 0{,}0021$, deviazione standard $\simeq 0{,}046$.}
}

\scheda{mlgauss}{
  \usoquando{Si stima la media di misure con rumore gaussiano di varianza nota.}
  \usocome{Si annulla la derivata della log-verosimiglianza: la stima è la media campione.}
  \usoesempio{25 misure con $\sigma = 1$: varianza della stima $\frac{1}{25} = 0{,}04$, deviazione standard $0{,}2$,
    uguale al limite di Cramér-Rao.}
}
